\section{Discussion}
\label{sec:discussion}

\para{Why does one iteration suffice?}
The first critique identifies the most salient weaknesses---predictable plot choices, underdeveloped characters, missed emotional beats---and the revision addresses them directly.
By iteration 2, the remaining issues are subtler, and the critique provides less actionable feedback.
By iteration 3, we observe signs of overcorrection: coherence declines ($-0.20$) as the model introduces plot elements to satisfy creative critique at the expense of narrative consistency.
This pattern aligns with the diminishing returns observed in Self-Refine across other tasks~\citep{madaan2023selfrefine} and with concerns about iterative self-refinement diverging from human quality judgments~\citep{pan2024rewardhacking}.

\para{Why does self-selection fail?}
The failure of \bestofn without refinement is striking and practically important.
The model generates four stories with genuine diversity, yet cannot reliably identify the best one.
This suggests that \emph{critique} (identifying specific weaknesses) and \emph{ranking} (holistic quality comparison) are distinct capabilities, and that \gptfour is substantially better at the former.
When \bestofn is combined with one refinement pass, performance recovers to 87\% win rate---comparable to \selfrefine $\times$2---indicating that starting from a better initial story (via selection) and then refining it is a valid alternative strategy.

\para{What does self-critique unlock?}
The dimension-level analysis (\secref{sec:dimensions}) reveals that self-critique primarily improves creative dimensions: originality, character depth, and emotional resonance.
Coherence, which is already at ceiling, does not benefit.
This pattern suggests that LLMs possess latent creative capacity---knowledge of what would make a more interesting, surprising, or emotionally affecting story---that they do not express in first drafts.
Default generation tends toward ``safe'' choices; structured self-reflection provides the model an opportunity to recognize and correct this conservatism.
This interpretation aligns with findings from \citet{jia2025writingzero}, who observed that self-critique-based training pushes models toward more creative outputs.

\para{Comparison with multi-agent approaches.}
\citet{bae2024critics} argued that single-agent critique is insufficient for creative improvement, advocating for multi-agent systems with diverse critic personas.
Our results partially challenge this claim: single-agent self-critique with a structured rubric produces very large improvements ($d = 2.02$) in a controlled blind evaluation.
However, the two setups are not directly comparable---\textsc{CritiCS} targets longer, more complex narratives and uses a different evaluation protocol.
Whether multi-agent critique offers additional gains \emph{on top of} structured single-agent critique for short fiction remains an open question.

\subsection{Limitations}
\label{sec:limitations}

\para{Same model for all roles.}
We use \gptfour for generation, critique, and evaluation.
While we employ different temperatures, seeds, and separate contexts for each role, shared model biases may inflate perceived improvement.
The model may systematically favor the kinds of revisions it tends to make, creating a feedback loop that does not reflect genuine quality gains.
Human evaluation on a subset of stories would substantially strengthen these findings.

\para{Single model family.}
Our results are specific to \gptfour and may not generalize to other architectures.
Open-source models and smaller models may lack the capacity for effective self-critique, as the critic must be sophisticated enough to identify genuine weaknesses.

\para{Inference-time only.}
We test self-critique at inference time; the model does not permanently learn from its critiques.
Training-time approaches~\citep{jia2025writingzero, yuan2024selfrewarding} may yield larger or more durable improvements.

\para{Scale of evaluation.}
Our 15 prompts are sufficient for detecting the large effects we observe but provide limited statistical power for subtler analyses (\eg comparing \selfrefine $\times$1 against $\times$2).
A larger prompt set would yield tighter confidence intervals.

\para{Short fiction only.}
We test stories of 400--600 words.
Self-critique effectiveness for longer narratives---where structural coherence becomes more challenging and single-pass revision may be insufficient---is untested.
